% --- Archon physics formalization source begin ---
% archon:physics
% archon:covers PhyXMiniProblems/problem_phyx_mini_0352.lean
% archon:source-report reports/phyx_mini/problem_phyx_mini_0352.source.json

\chapter{Physics problem phyx\_mini\_0352}
\label{ch:PhyXMiniProblems_problem_phyx_mini_0352}

\paragraph{Problem source.}
\#\# Physical scenario

A monatomic ideal gas is taken around the cycle shown in figure in the direction shown in the figure. The path for process \$c \textbackslash{}rightarrow a\$ is a straight line in the \$pV\$-diagram.

\#\# Question

What is the efficiency of the cycle?

\#\# Answer choices

- A: A: 10\textbackslash{}\%

- B: B: 11.1\textbackslash{}\%

- C: C: 15.9\textbackslash{}\%

- D: D: 12.3\textbackslash{}\%

\#\# Figure caption (auxiliary; use the image as primary evidence)

The image is a pressure-volume (P-V) diagram of a thermodynamic process. It features the following components:

- **Axes**: 
  - The vertical axis is labeled “\textbackslash{}(p\textbackslash{})” for pressure, with markings at \textbackslash{}(1.00 \textbackslash{}times 10\textasciicircum{}5 \textbackslash{}, \textbackslash{}text\{Pa\}\textbackslash{}) and \textbackslash{}(3.00 \textbackslash{}times 10\textasciicircum{}5 \textbackslash{}, \textbackslash{}text\{Pa\}\textbackslash{}).
  - The horizontal axis is labeled “\textbackslash{}(V\textbackslash{})” for volume, with markings at \textbackslash{}(0.500\textbackslash{}, \textbackslash{}text\{m\}\textasciicircum{}3\textbackslash{}) and \textbackslash{}(0.800\textbackslash{}, \textbackslash{}text\{m\}\textasciicircum{}3\textbackslash{}).

- **Points**:
  - Point \textbackslash{}(a\textbackslash{}) is at the coordinate \textbackslash{}((0.500\textbackslash{}, \textbackslash{}text\{m\}\textasciicircum{}3, 3.00 \textbackslash{}times 10\textasciicircum{}5 \textbackslash{}, \textbackslash{}text\{Pa\})\textbackslash{}).
  - Point \textbackslash{}(b\textbackslash{}) is at \textbackslash{}((0.800\textbackslash{}, \textbackslash{}text\{m\}\textasciicircum{}3, 3.00 \textbackslash{}times 10\textasciicircum{}5 \textbackslash{}, \textbackslash{}text\{Pa\})\textbackslash{}).
  - Point \textbackslash{}(c\textbackslash{}) is at \textbackslash{}((0.800\textbackslash{}, \textbackslash{}text\{m\}\textasciicircum{}3, 1.00 \textbackslash{}times 10\textasciicircum{}5 \textbackslash{}, \textbackslash{}text\{Pa\})\textbackslash{}).

- **Process Arrows**:
  - An arrow from \textbackslash{}(a\textbackslash{}) to \textbackslash{}(b\textbackslash{}) indicates an isobaric expansion.
  - An arrow from \textbackslash{}(b\textbackslash{}) to \textbackslash{}(c\textbackslash{}) indicates an isochoric decrease in pressure.
  - An arrow from \textbackslash{}(c\textbackslash{}) to \textbackslash{}(a\textbackslash{}) indicates an isochoric increase in pressure, completing a cyclic process.

The diagram depicts a cycle involving processes at constant pressure and constant volume.

\#\# Dataset metadata

Laws of Thermodynamics; Multi-Formula Reasoning, Physical Model Grounding Reasoning

\paragraph{Recorded answer/context.}
B: B: 11.1\textbackslash{}\%

\paragraph{Figure/image path.}
/root/proposal\_for\_physic/science-mango/phyx\_mini\_run/phyx\_data/test\_image/352.png

\paragraph{Formalization target.}
create a compiling Lean file with sorry bodies at `PhyXMiniProblems/problem\_phyx\_mini\_0352.lean`. The Lean declarations must preserve the physical quantities, dimensions or dimensional roles, figure labels, governing-law hypotheses, and final relation expressed by this problem.
Use Mathlib/Physlib names found through LeanExplore where available. If a physics API is missing, introduce faithful local abstractions rather than scalar placeholder aliases.

\begin{theorem}[Physics formalization target]
\label{thm:physics:phyx_mini_0352:target}
\lean{PhyXMiniProblems.ProblemPhyXMini0352.problem_phyx_mini_0352}
\uses{def:physics:phyx-mini-0352:phyxminiproblems-problemphyxmini0352-monatomicidealgascycle, def:physics:phyx-mini-0352:phyxminiproblems-problemphyxmini0352-matchesproblemandprimaryfigure, def:physics:phyx-mini-0352:phyxminiproblems-problemphyxmini0352-hasphysicalcycleparameters, def:physics:phyx-mini-0352:phyxminiproblems-problemphyxmini0352-satisfiesmonatomicidealgasinternalenergylaw, def:physics:phyx-mini-0352:phyxminiproblems-problemphyxmini0352-satisfiesstraightpathboundaryworklaw, def:physics:phyx-mini-0352:phyxminiproblems-problemphyxmini0352-satisfiesfirstlawoneachleg, def:physics:phyx-mini-0352:phyxminiproblems-problemphyxmini0352-satisfiesheatengineefficiencylaw, def:physics:phyx-mini-0352:phyxminiproblems-problemphyxmini0352-recordedanswerchoice, def:physics:phyx-mini-0352:phyxminiproblems-problemphyxmini0352-efficiencymatchesdisplayedchoice}
The assigned autoformalize agent should translate this physics problem into Lean declarations in the covered file, with theorem and lemma proof bodies written as `by sorry`.
\end{theorem}
\begin{proof}
This is an autoformalization task, not a proof task. Produce statements that can later be proved by the physics prover without weakening the physical model.
\end{proof}

% --- Archon declaration topology begin ---
\section{Declaration topology}
\begin{definition}[Volume Quantity]
  \label{def:physics:phyx-mini-0352:phyxminiproblems-problemphyxmini0352-volumequantity}
  \lean{PhyXMiniProblems.ProblemPhyXMini0352.VolumeQuantity}
  A signed physical volume, carrying the dimension of length cubed
\end{definition}
\begin{proof}
  This is the declared model definition of Volume Quantity.
\end{proof}
\begin{definition}[Pressure Quantity]
  \label{def:physics:phyx-mini-0352:phyxminiproblems-problemphyxmini0352-pressurequantity}
  \lean{PhyXMiniProblems.ProblemPhyXMini0352.PressureQuantity}
  A physical pressure, using Physlib's pressure dimension
\end{definition}
\begin{proof}
  This is the declared model definition of Pressure Quantity.
\end{proof}
\begin{definition}[Energy Quantity]
  \label{def:physics:phyx-mini-0352:phyxminiproblems-problemphyxmini0352-energyquantity}
  \lean{PhyXMiniProblems.ProblemPhyXMini0352.EnergyQuantity}
  A signed physical energy, used for internal energy, heat, and work
\end{definition}
\begin{proof}
  This is the declared model definition of Energy Quantity.
\end{proof}
\begin{definition}[volume In Cubic Meters]
  \label{def:physics:phyx-mini-0352:phyxminiproblems-problemphyxmini0352-volumeincubicmeters}
  \lean{PhyXMiniProblems.ProblemPhyXMini0352.volumeInCubicMeters}
  \uses{def:physics:phyx-mini-0352:phyxminiproblems-problemphyxmini0352-volumequantity}
  SI readout of a volume, in cubic metres
\end{definition}
\begin{proof}
  This is the declared model definition of volume In Cubic Meters.
\end{proof}
\begin{definition}[pressure In Pascals]
  \label{def:physics:phyx-mini-0352:phyxminiproblems-problemphyxmini0352-pressureinpascals}
  \lean{PhyXMiniProblems.ProblemPhyXMini0352.pressureInPascals}
  \uses{def:physics:phyx-mini-0352:phyxminiproblems-problemphyxmini0352-pressurequantity}
  SI readout of a pressure, in pascals
\end{definition}
\begin{proof}
  This is the declared model definition of pressure In Pascals.
\end{proof}
\begin{definition}[energy In Joules]
  \label{def:physics:phyx-mini-0352:phyxminiproblems-problemphyxmini0352-energyinjoules}
  \lean{PhyXMiniProblems.ProblemPhyXMini0352.energyInJoules}
  \uses{def:physics:phyx-mini-0352:phyxminiproblems-problemphyxmini0352-energyquantity}
  SI readout of a signed energy, in joules
\end{definition}
\begin{proof}
  This is the declared model definition of energy In Joules.
\end{proof}
\begin{definition}[Gas Model]
  \label{def:physics:phyx-mini-0352:phyxminiproblems-problemphyxmini0352-gasmodel}
  \lean{PhyXMiniProblems.ProblemPhyXMini0352.GasModel}
  The material model named in the problem statement
\end{definition}
\begin{proof}
  This is the declared model definition of Gas Model.
\end{proof}
\begin{definition}[Cycle State]
  \label{def:physics:phyx-mini-0352:phyxminiproblems-problemphyxmini0352-cyclestate}
  \lean{PhyXMiniProblems.ProblemPhyXMini0352.CycleState}
  The three labeled equilibrium states in the pV diagram
\end{definition}
\begin{proof}
  This is the declared model definition of Cycle State.
\end{proof}
\begin{definition}[Cycle Leg]
  \label{def:physics:phyx-mini-0352:phyxminiproblems-problemphyxmini0352-cycleleg}
  \lean{PhyXMiniProblems.ProblemPhyXMini0352.CycleLeg}
  The directed legs of the cycle, in traversal order
\end{definition}
\begin{proof}
  This is the declared model definition of Cycle Leg.
\end{proof}
\begin{definition}[Process Kind]
  \label{def:physics:phyx-mini-0352:phyxminiproblems-problemphyxmini0352-processkind}
  \lean{PhyXMiniProblems.ProblemPhyXMini0352.ProcessKind}
  Thermodynamic character singled out for each process leg
\end{definition}
\begin{proof}
  This is the declared model definition of Process Kind.
\end{proof}
\begin{definition}[Path Shape]
  \label{def:physics:phyx-mini-0352:phyxminiproblems-problemphyxmini0352-pathshape}
  \lean{PhyXMiniProblems.ProblemPhyXMini0352.PathShape}
  Geometric shape of a process in the pressure-volume plane
\end{definition}
\begin{proof}
  This is the declared model definition of Path Shape.
\end{proof}
\begin{definition}[Arrow Direction]
  \label{def:physics:phyx-mini-0352:phyxminiproblems-problemphyxmini0352-arrowdirection}
  \lean{PhyXMiniProblems.ProblemPhyXMini0352.ArrowDirection}
  Directions of the three blue arrows in the primary image
\end{definition}
\begin{proof}
  This is the declared model definition of Arrow Direction.
\end{proof}
\begin{definition}[Axis Label]
  \label{def:physics:phyx-mini-0352:phyxminiproblems-problemphyxmini0352-axislabel}
  \lean{PhyXMiniProblems.ProblemPhyXMini0352.AxisLabel}
  Literal mathematical labels shown on the axes and at the origin
\end{definition}
\begin{proof}
  This is the declared model definition of Axis Label.
\end{proof}
\begin{definition}[Monatomic Ideal Gas Cycle]
  \label{def:physics:phyx-mini-0352:phyxminiproblems-problemphyxmini0352-monatomicidealgascycle}
  \lean{PhyXMiniProblems.ProblemPhyXMini0352.MonatomicIdealGasCycle}
  \uses{def:physics:phyx-mini-0352:phyxminiproblems-problemphyxmini0352-volumequantity, def:physics:phyx-mini-0352:phyxminiproblems-problemphyxmini0352-pressurequantity, def:physics:phyx-mini-0352:phyxminiproblems-problemphyxmini0352-energyquantity, def:physics:phyx-mini-0352:phyxminiproblems-problemphyxmini0352-gasmodel, def:physics:phyx-mini-0352:phyxminiproblems-problemphyxmini0352-cyclestate, def:physics:phyx-mini-0352:phyxminiproblems-problemphyxmini0352-cycleleg, def:physics:phyx-mini-0352:phyxminiproblems-problemphyxmini0352-processkind, def:physics:phyx-mini-0352:phyxminiproblems-problemphyxmini0352-pathshape, def:physics:phyx-mini-0352:phyxminiproblems-problemphyxmini0352-arrowdirection, def:physics:phyx-mini-0352:phyxminiproblems-problemphyxmini0352-axislabel}
  The independent physical observables of the gas cycle. The efficiency is retained as an independent dimensionless observable and is constrained only by the general efficiency law below. In particular, it is not defined to be the requested numerical answer
\end{definition}
\begin{proof}
  This is the declared model definition of Monatomic Ideal Gas Cycle.
\end{proof}
\begin{definition}[cycle Net Work In Joules]
  \label{def:physics:phyx-mini-0352:phyxminiproblems-problemphyxmini0352-cyclenetworkinjoules}
  \lean{PhyXMiniProblems.ProblemPhyXMini0352.cycleNetWorkInJoules}
  \uses{def:physics:phyx-mini-0352:phyxminiproblems-problemphyxmini0352-energyinjoules, def:physics:phyx-mini-0352:phyxminiproblems-problemphyxmini0352-monatomicidealgascycle}
  Total work done by the gas over one traversal of the three-leg cycle
\end{definition}
\begin{proof}
  This is the declared model definition of cycle Net Work In Joules.
\end{proof}
\begin{definition}[positive Heat In Joules]
  \label{def:physics:phyx-mini-0352:phyxminiproblems-problemphyxmini0352-positiveheatinjoules}
  \lean{PhyXMiniProblems.ProblemPhyXMini0352.positiveHeatInJoules}
  \uses{def:physics:phyx-mini-0352:phyxminiproblems-problemphyxmini0352-energyquantity, def:physics:phyx-mini-0352:phyxminiproblems-problemphyxmini0352-energyinjoules}
  The positive part of a signed heat transfer into the gas
\end{definition}
\begin{proof}
  This is the declared model definition of positive Heat In Joules.
\end{proof}
\begin{definition}[cycle Heat Input In Joules]
  \label{def:physics:phyx-mini-0352:phyxminiproblems-problemphyxmini0352-cycleheatinputinjoules}
  \lean{PhyXMiniProblems.ProblemPhyXMini0352.cycleHeatInputInJoules}
  \uses{def:physics:phyx-mini-0352:phyxminiproblems-problemphyxmini0352-monatomicidealgascycle, def:physics:phyx-mini-0352:phyxminiproblems-problemphyxmini0352-positiveheatinjoules}
  Total heat absorbed by the gas, excluding legs on which heat leaves it
\end{definition}
\begin{proof}
  This is the declared model definition of cycle Heat Input In Joules.
\end{proof}
\begin{definition}[Matches Problem And Primary Figure]
  \label{def:physics:phyx-mini-0352:phyxminiproblems-problemphyxmini0352-matchesproblemandprimaryfigure}
  \lean{PhyXMiniProblems.ProblemPhyXMini0352.MatchesProblemAndPrimaryFigure}
  \uses{def:physics:phyx-mini-0352:phyxminiproblems-problemphyxmini0352-volumeincubicmeters, def:physics:phyx-mini-0352:phyxminiproblems-problemphyxmini0352-pressureinpascals, def:physics:phyx-mini-0352:phyxminiproblems-problemphyxmini0352-monatomicidealgascycle}
  Exact transcription of the problem statement and the primary raster image. The state coordinates are physical SI readouts. The path endpoints and arrow directions record the displayed clockwise traversal, while the process tags record the horizontal, vertical, and diagonal legs. No requested efficiency or answer-choice value occurs in this structure
\end{definition}
\begin{proof}
  This is the declared model definition of Matches Problem And Primary Figure.
\end{proof}
\begin{definition}[Has Physical Cycle Parameters]
  \label{def:physics:phyx-mini-0352:phyxminiproblems-problemphyxmini0352-hasphysicalcycleparameters}
  \lean{PhyXMiniProblems.ProblemPhyXMini0352.HasPhysicalCycleParameters}
  \uses{def:physics:phyx-mini-0352:phyxminiproblems-problemphyxmini0352-volumeincubicmeters, def:physics:phyx-mini-0352:phyxminiproblems-problemphyxmini0352-pressureinpascals, def:physics:phyx-mini-0352:phyxminiproblems-problemphyxmini0352-monatomicidealgascycle, def:physics:phyx-mini-0352:phyxminiproblems-problemphyxmini0352-cycleheatinputinjoules}
  Positivity and nondegeneracy conditions for a physical heat-engine cycle
\end{definition}
\begin{proof}
  This is the declared model definition of Has Physical Cycle Parameters.
\end{proof}
\begin{definition}[Satisfies Monatomic Ideal Gas Internal Energy Law]
  \label{def:physics:phyx-mini-0352:phyxminiproblems-problemphyxmini0352-satisfiesmonatomicidealgasinternalenergylaw}
  \lean{PhyXMiniProblems.ProblemPhyXMini0352.SatisfiesMonatomicIdealGasInternalEnergyLaw}
  \uses{def:physics:phyx-mini-0352:phyxminiproblems-problemphyxmini0352-volumeincubicmeters, def:physics:phyx-mini-0352:phyxminiproblems-problemphyxmini0352-pressureinpascals, def:physics:phyx-mini-0352:phyxminiproblems-problemphyxmini0352-energyinjoules, def:physics:phyx-mini-0352:phyxminiproblems-problemphyxmini0352-monatomicidealgascycle}
  For a monatomic ideal gas, U = (3/2) nRT = (3/2) pV at every equilibrium state. This is the only ideal-gas caloric relation needed for the cycle; it contains no figure-specific coordinates or requested efficiency
\end{definition}
\begin{proof}
  This is the declared model definition of Satisfies Monatomic Ideal Gas Internal Energy Law.
\end{proof}
\begin{definition}[Satisfies Straight Path Boundary Work Law]
  \label{def:physics:phyx-mini-0352:phyxminiproblems-problemphyxmini0352-satisfiesstraightpathboundaryworklaw}
  \lean{PhyXMiniProblems.ProblemPhyXMini0352.SatisfiesStraightPathBoundaryWorkLaw}
  \uses{def:physics:phyx-mini-0352:phyxminiproblems-problemphyxmini0352-volumeincubicmeters, def:physics:phyx-mini-0352:phyxminiproblems-problemphyxmini0352-pressureinpascals, def:physics:phyx-mini-0352:phyxminiproblems-problemphyxmini0352-energyinjoules, def:physics:phyx-mini-0352:phyxminiproblems-problemphyxmini0352-monatomicidealgascycle}
  Boundary work for any straight segment in the pV plane. Linear pressure variation makes the work by the gas equal to average pressure times the change in volume. The formula covers the isobaric, isochoric, and diagonal legs uniformly and is independent of this problem's numerical data
\end{definition}
\begin{proof}
  This is the declared model definition of Satisfies Straight Path Boundary Work Law.
\end{proof}
\begin{definition}[Satisfies First Law On Each Leg]
  \label{def:physics:phyx-mini-0352:phyxminiproblems-problemphyxmini0352-satisfiesfirstlawoneachleg}
  \lean{PhyXMiniProblems.ProblemPhyXMini0352.SatisfiesFirstLawOnEachLeg}
  \uses{def:physics:phyx-mini-0352:phyxminiproblems-problemphyxmini0352-energyinjoules, def:physics:phyx-mini-0352:phyxminiproblems-problemphyxmini0352-monatomicidealgascycle}
  First law on each directed leg, using heat into the gas and work by the gas as positive: Q = ΔU + W. Heat transfers are independent observables rather than definitions chosen to produce the final efficiency
\end{definition}
\begin{proof}
  This is the declared model definition of Satisfies First Law On Each Leg.
\end{proof}
\begin{definition}[Satisfies Heat Engine Efficiency Law]
  \label{def:physics:phyx-mini-0352:phyxminiproblems-problemphyxmini0352-satisfiesheatengineefficiencylaw}
  \lean{PhyXMiniProblems.ProblemPhyXMini0352.SatisfiesHeatEngineEfficiencyLaw}
  \uses{def:physics:phyx-mini-0352:phyxminiproblems-problemphyxmini0352-monatomicidealgascycle, def:physics:phyx-mini-0352:phyxminiproblems-problemphyxmini0352-cyclenetworkinjoules, def:physics:phyx-mini-0352:phyxminiproblems-problemphyxmini0352-cycleheatinputinjoules}
  Definition of thermal efficiency for a cyclic heat engine: net work output divided by total positive heat input. This general relation does not specify the efficiency of the present numerical cycle
\end{definition}
\begin{proof}
  This is the declared model definition of Satisfies Heat Engine Efficiency Law.
\end{proof}
\begin{lemma}[cycle energy accounting]
  \label{lem:physics:phyx-mini-0352:phyxminiproblems-problemphyxmini0352-cycle-energy-accounting}
  \lean{PhyXMiniProblems.ProblemPhyXMini0352.cycle_energy_accounting}
  \uses{def:physics:phyx-mini-0352:phyxminiproblems-problemphyxmini0352-energyinjoules, def:physics:phyx-mini-0352:phyxminiproblems-problemphyxmini0352-monatomicidealgascycle, def:physics:phyx-mini-0352:phyxminiproblems-problemphyxmini0352-cyclenetworkinjoules, def:physics:phyx-mini-0352:phyxminiproblems-problemphyxmini0352-cycleheatinputinjoules, def:physics:phyx-mini-0352:phyxminiproblems-problemphyxmini0352-matchesproblemandprimaryfigure, def:physics:phyx-mini-0352:phyxminiproblems-problemphyxmini0352-satisfiesmonatomicidealgasinternalenergylaw, def:physics:phyx-mini-0352:phyxminiproblems-problemphyxmini0352-satisfiesstraightpathboundaryworklaw, def:physics:phyx-mini-0352:phyxminiproblems-problemphyxmini0352-satisfiesfirstlawoneachleg}
  The thermodynamic laws and figure readouts determine all three signed works, state internal energies, and signed heats. Positive work and heat are into the output/work and gas/heat conventions described above
\end{lemma}
\begin{proof}
  The conclusion follows from the governing relations and figure readouts named in the statement of cycle energy accounting.
\end{proof}
\begin{definition}[Answer Choice]
  \label{def:physics:phyx-mini-0352:phyxminiproblems-problemphyxmini0352-answerchoice}
  \lean{PhyXMiniProblems.ProblemPhyXMini0352.AnswerChoice}
  Labels of the four multiple-choice answers
\end{definition}
\begin{proof}
  This is the declared model definition of Answer Choice.
\end{proof}
\begin{definition}[displayed Efficiency Percent]
  \label{def:physics:phyx-mini-0352:phyxminiproblems-problemphyxmini0352-displayedefficiencypercent}
  \lean{PhyXMiniProblems.ProblemPhyXMini0352.displayedEfficiencyPercent}
  \uses{def:physics:phyx-mini-0352:phyxminiproblems-problemphyxmini0352-answerchoice}
  Percentage printed beside each answer label
\end{definition}
\begin{proof}
  This is the declared model definition of displayed Efficiency Percent.
\end{proof}
\begin{definition}[recorded Answer Choice]
  \label{def:physics:phyx-mini-0352:phyxminiproblems-problemphyxmini0352-recordedanswerchoice}
  \lean{PhyXMiniProblems.ProblemPhyXMini0352.recordedAnswerChoice}
  \uses{def:physics:phyx-mini-0352:phyxminiproblems-problemphyxmini0352-answerchoice}
  The answer label recorded in the supplied dataset metadata
\end{definition}
\begin{proof}
  This is the declared model definition of recorded Answer Choice.
\end{proof}
\begin{definition}[efficiency Percent]
  \label{def:physics:phyx-mini-0352:phyxminiproblems-problemphyxmini0352-efficiencypercent}
  \lean{PhyXMiniProblems.ProblemPhyXMini0352.efficiencyPercent}
  Convert a dimensionless efficiency to a percentage readout
\end{definition}
\begin{proof}
  This is the declared model definition of efficiency Percent.
\end{proof}
\begin{definition}[Efficiency Matches Displayed Choice]
  \label{def:physics:phyx-mini-0352:phyxminiproblems-problemphyxmini0352-efficiencymatchesdisplayedchoice}
  \lean{PhyXMiniProblems.ProblemPhyXMini0352.EfficiencyMatchesDisplayedChoice}
  \uses{def:physics:phyx-mini-0352:phyxminiproblems-problemphyxmini0352-answerchoice, def:physics:phyx-mini-0352:phyxminiproblems-problemphyxmini0352-displayedefficiencypercent, def:physics:phyx-mini-0352:phyxminiproblems-problemphyxmini0352-efficiencypercent}
  Agreement with a percentage displayed to the nearest tenth of a percent
\end{definition}
\begin{proof}
  This is the declared model definition of Efficiency Matches Displayed Choice.
\end{proof}
% --- Archon declaration topology end ---
% --- Archon physics formalization source end ---
